
\chapter{Complessit\`a}

\section{P, NP e NP-completo}
I problemi P, NP e NP-completi sono concetti fondamentali nella teoria della complessità computazionale. 

Il problema P rappresenta la classe dei problemi che possono essere risolti in tempo polinomiale da un algoritmo deterministico. In altre parole, per un problema appartenente a P, esiste un algoritmo che può risolverlo in un tempo che cresce al più in modo polinomiale rispetto alla dimensione dell'input.

Il problema NP rappresenta la classe dei problemi per i quali una soluzione proposta può essere verificata in tempo polinomiale da un algoritmo deterministico. In altre parole, se si ha una soluzione proposta per un problema NP, è possibile verificarne la correttezza in un tempo che cresce al più in modo polinomiale rispetto alla dimensione dell'input.

I problemi NP-completi sono una sottoclasse dei problemi NP che sono considerati tra i più difficili da risolvere. Un problema NP-completo è un problema per il quale ogni altro problema in NP può essere ridotto in modo polinomiale ad esso. In altre parole, se si trovasse un algoritmo efficiente per risolvere un problema NP-completo, si potrebbe risolvere in modo efficiente tutti i problemi in NP.

Tuttavia, non si sa ancora se P è uguale a NP o se P è diverso da NP. Questa è una delle domande aperte più importanti nella teoria della complessità computazionale. La maggior parte degli esperti ritiene che P sia diverso da NP, ma finora non è stata trovata una prova definitiva.
